\begin{lemma}
For every \(\varepsilon>0\), there are \(\delta>0\) and \(s_0\) such that whenever \(s\ge s_0\) and \(a\ge 0\) is an integer satisfying \(a\le \delta s\), there exists a loopless \(T_s\)-free digraph \(D\) on
\[
  2^{2s-3}-2^{s-1}-2^{s-2}+1
\]
vertices such that
\[
  \overrightarrow{i_{s+a}}(D)\le
  2^{\frac32s^2+as-\frac52s+\varepsilon s}.
\]
\end{lemma}
\begin{proof}
Fix \(\varepsilon>0\). All logarithms in this proof are base \(2\). Since
\[
  x\log_2\left(\frac{2e}{x}\right)\to0
  \qquad\text{as }x\to0^+,
\]
we may choose \(0<\delta\le1\) such that
\[
  x\log_2\left(\frac{2e}{x}\right)\le \frac{\varepsilon}{4}
  \qquad\text{for every }0<x\le\delta. \tag{1}
\]
Choose \(s_0\ge4\) large enough that, for every \(s\ge s_0\),
\[
  \frac12\left(\log_2(2es)+\frac32\right)^2\le \frac{\varepsilon s}{4},
  \qquad
  \log_2 s\le\frac{\varepsilon s}{2}. \tag{2}
\]
This is possible because \((\log s)^2=o(s)\) and \(\log s=o(s)\).

Now let \(s\ge s_0\), let \(a\ge0\) be an integer with \(a\le\delta s\), and set
\[
  k=s+a.
\]
Then \(4\le s\le k\), so \(\noderef{F2ForwardIndependentTuples}\) applies. It gives a finite vertex set and a digraph \(D\) which is loopless and \(T_s\)-free, has exactly
\[
  2^{2s-3}-2^{s-1}-2^{s-2}+1
\]
vertices, and satisfies
\[
  \overrightarrow{i_k}(D)\le
  \sum_{t=1}^{s-1}
    \binom{k}{t}2^{(s-1)(t+k)-\binom{t+1}{2}}. \tag{3}
\]
Since \(k=s+a\), this is the desired digraph once we prove that the sum in (3) is at most
\[
  2^{\frac32s^2+as-\frac52s+\varepsilon s}. \tag{4}
\]

For \(1\le t\le s-1\), denote the \(t\)-th summand in (3) by
\[
  M_t=\binom{k}{t}2^{(s-1)(t+k)-\binom{t+1}{2}}.
\]
Write \(t=s-j\). Then \(1\le j\le s-1\), and
\[
  \binom{k}{t}
  =\binom{s+a}{s-j}
  =\binom{s+a}{j+a}.
\]
Using the standard estimate \(\binom{n}{m}\le (en/m)^m\), and using \(j+a\ge a+1\), we get
\[
  \binom{s+a}{j+a}
  \le
  \left(\frac{e(s+a)}{j+a}\right)^{j+a}
  \le
  \left(\frac{e(s+a)}{a+1}\right)^{j+a}. \tag{5}
\]
The exponent coming from the power of \(2\) is
\[
\begin{aligned}
  (s-1)(s-j+k)-\binom{s-j+1}{2}
  &=
  \frac{s^2}{2}+sk-\frac{3s}{2}-k+\frac{3j}{2}-\frac{j^2}{2}.
\end{aligned} \tag{6}
\]
Combining (5) and (6), and putting
\[
  R=\frac{e(s+a)}{a+1},\qquad \lambda=\log_2 R,
\]
we obtain
\[
  M_{s-j}
  \le
  2^{
    \frac{s^2}{2}+sk-\frac{3s}{2}-k
    +a\lambda+j\lambda+\frac{3j}{2}-\frac{j^2}{2}
  }. \tag{7}
\]

We next bound the two extra terms in (7), uniformly in \(0\le a\le\delta s\). First consider \(a\lambda\). If \(a=0\), this term is \(0\). If \(a>0\), put \(x=a/s\). Then \(0<x\le\delta\), and since \(a+1\ge a=xs\) and \(s+a\le(1+\delta)s\le2s\),
\[
  R=\frac{e(s+a)}{a+1}
  \le \frac{2e}{x}.
\]
Therefore, by (1),
\[
  a\lambda
  =xs\log_2 R
  \le xs\log_2\left(\frac{2e}{x}\right)
  \le \frac{\varepsilon s}{4}. \tag{8}
\]
Thus (8) also holds when \(a=0\).

Now bound the \(j\)-dependent part. Since \(a\le\delta s\) and \(\delta\le1\), we have
\[
  R=\frac{e(s+a)}{a+1}\le e(s+a)\le2es,
\]
so \(\lambda\le\log_2(2es)\). For real \(y\ge0\), the quadratic
\[
  y\lambda+\frac{3y}{2}-\frac{y^2}{2}
\]
is maximized at \(y=\lambda+\frac32\), and its maximum is
\[
  \frac12\left(\lambda+\frac32\right)^2.
\]
Since \(j\ge0\), (2) gives
\[
  j\lambda+\frac{3j}{2}-\frac{j^2}{2}
  \le
  \frac12\left(\log_2(2es)+\frac32\right)^2
  \le \frac{\varepsilon s}{4}. \tag{9}
\]

Substituting (8) and (9) into (7), every summand in (3) satisfies
\[
  M_t\le
  2^{\frac{s^2}{2}+sk-\frac{3s}{2}-k+\frac{\varepsilon s}{2}}
  \qquad(1\le t\le s-1). \tag{10}
\]
There are at most \(s\) summands, and the second inequality in (2) gives
\[
\begin{aligned}
  \sum_{t=1}^{s-1}M_t
  &\le
  s\cdot
  2^{\frac{s^2}{2}+sk-\frac{3s}{2}-k+\frac{\varepsilon s}{2}}  \\
  &=
  2^{\frac{s^2}{2}+sk-\frac{3s}{2}-k+\frac{\varepsilon s}{2}+\log_2 s} \\
  &\le
  2^{\frac{s^2}{2}+sk-\frac{3s}{2}-k+\varepsilon s}.
\end{aligned} \tag{11}
\]
Finally substitute \(k=s+a\). The exponent in (11) becomes
\[
  \frac{s^2}{2}+s(s+a)-\frac{3s}{2}-(s+a)+\varepsilon s
  =
  \frac32s^2+as-\frac52s-a+\varepsilon s
  \le
  \frac32s^2+as-\frac52s+\varepsilon s,
\]
because \(a\ge0\). Combining this with (3) proves (4), and hence
\[
  \overrightarrow{i_{s+a}}(D)\le
  2^{\frac32s^2+as-\frac52s+\varepsilon s}.
\]
The same application of \(\noderef{F2ForwardIndependentTuples}\) already supplied looplessness, \(T_s\)-freeness, and the exact vertex count, so the lemma follows.
\end{proof}
